% P16-C4: UAI 2006 IDC, exact exchange closure and conditional soundness.
\section{Conditional interventional identification}
\label{sec:p16-conditional-interventional-identification}

Shpitser and Pearl's UAI 2006 procedure first exchanges eligible
observations for actions, and only then asks the joint identification
algorithm for the distribution to normalize
\cite[Section~5, Figure~3 and Appendix]{ShpitserPearl2006ConditionalIdentification}.
The author-hosted UAI version and its separate 2007 addendum are distinct
sources. The addendum retracts the original paper's assertion that Tian's
conditional identification algorithm is unsound; that assertion is not
adopted here \cite{Shpitser2007ConditionalIdentificationAppendum}.

\begin{hbdefinition}{def:p16-idc-query}{Conditional effects and their quantified model class}
Fix the finite independent-latent-input class on an observed mixed graph
\(G\), with observed vertices \(N\) and finite nonempty value sets
\(\mathbf S=(S_i)_{i\in N}\). Let \(X,Y,Z\subseteq N\)
be pairwise disjoint, with assignments \(x,y,z\). In a specified
model define
\[
 P_x(y\mid z)=\frac{P_x(y,z)}{P_x(z)}
 \qquad\text{when }P_x(z)>0.
\]
The evidence belongs to the world after \(X=x\); it is not
an observation in a separate counterfactual world.

The main execution domain here is a compatible full observed law
\(p(v)>0\) at every assignment. For one fixed \(\mathbf S\),
conditional identification means agreement of these kernels in every
two compatible positive-law models sharing any common observed law.
Graph-uniform conditional identification quantifies over every finite
nonempty \(\mathbf S\), every possible common positive law \(p\),
every such pair of models, and every \(x,y,z\).
The source separately permits nonpositive observed laws and requires
\(P_x(z)>0\); that further support domain is distinguished below.
\end{hbdefinition}

\begin{hblemma}{thm:p16-idc-positive-context}{proved=annex:p16.idc-positive-context}{Positive observations supply every conditional context in the main domain}
If \(p>0\) everywhere, then every intervention law is positive
at every assignment of its remaining observed variables. In particular,
all the denominators \(P_x(z)\) in the preceding definition and
all conditional contexts reached by IDC are positive.
\end{hblemma}

\begin{hbdefinition}{def:p16-idc-eligible}{An eligible observation and its graph license}
For graph tests let \(\widehat G\) be the canonical DAG obtained
by retaining the observed directed edges and replacing each bidirected
edge \(i\leftrightarrow j\) by one formal root with arrows to
\(i,j\). This is a separation device, not a replacement of a
model's actual latent distributions or child sets. For the current query \((X,Y,Z)\), an
observation vertex \(w\in Z\) is \emph{eligible} when
\[
 Y\mathrel{\perp_{\widehat G_{\overline X,\underline{\{w\}}}}}
       \{w\}\mid\bigl(X\cup(Z\setminus\{w\})\bigr).
\]
The overline deletes incoming arrows at the current action vertices;
the underline deletes outgoing arrows at the proposed exchange vertex.
Let \(E(G;X,Y,Z)\) be the set of eligible vertices. This is a
graphical set. An equality caused by one model's numerical choices
does not by itself make a vertex eligible.
\end{hbdefinition}

\begin{hblemma}{thm:p16-idc-root-graph-bridge}{proved=annex:p16.idc-root-graph-bridge}{The canonical separation test preserves actual latent-root separation}
For observed endpoint and conditioning sets, every separation test in
the canonical \(\widehat G\), including the incoming/outgoing
arrow deletions used by IDC, has the same result in the expansion
using any model's actual independent latent roots with observed graph
\(G\). No equivalence of the two statistical model classes is
asserted by this graph lemma.
\end{hblemma}

\begin{hblemma}{thm:p16-idc-rule-two-step}{proved=annex:p16.idc-rule-two-step}{The exact observation-to-action exchange}
For an eligible \(w\), at every supported current context,
\[
 P_x(y\mid z)=P_{x,z_w}(y\mid z_{Z\setminus\{w\}}).
\]
Positivity of the left conditioning event supplies the conditioning
support on the right by consistency. In the main positive-law domain
both sides exist for all assignments.
\end{hblemma}

\begin{hbtheorem}{thm:p16-idc-exchange-closure}{proved=annex:p16.idc-exchange-closure}{The unique maximal graphical exchange set}
Exchanging one eligible vertex leaves eligibility of every other
conditioning vertex unchanged. Consequently all orders of single-vertex
Rule 2 exchanges move exactly \(E=E(G;X,Y,Z)\), and terminate
with action set \(X\cup E\) and conditioning set \(Z\setminus E\).
The set \(E\) also satisfies the simultaneous Rule 2 graph
condition. Any set satisfying that simultaneous condition is a subset
of \(E\). Thus the maximal graphical exchange set is unique.
\end{hbtheorem}

\begin{hbmethod}{def:p16-idc-algorithm}{The IDC procedure}
Input \((Y,X,Z,p,G)\) has the preceding meanings; value assignments
remain explicit arguments of the returned function. Keep \(p,G\)
unchanged during these conditioning exchanges.
\begin{enumerate}
\item If \(Y=\varnothing\), return one. If an eligible
      \(w\in Z\) exists, call
      \[
       \mathsf{IDC}(Y,X\cup\{w\},Z\setminus\{w\},p,G),
      \]
      using the existing observed value \(z_w\) as its new
      action value. Retain this exchange and its graph license.
\item If no eligible vertex remains, call
      \(\mathsf{ID}(Y\cup Z,X,p,G)\). If ID succeeds with
      \(q_x(y,z)\), return
      \[
       \frac{q_x(y,z)}{\sum_{y'\in S_Y}q_x(y',z)}.
      \]
      The sum holds the action and conditioning assignments fixed.
      If ID fails, retain its full failure record together with
      the preceding exchange sequence; do not divide a failure.
\end{enumerate}
When \(Z\) is empty the terminal denominator is one. The
empty-response return makes explicit a case already implied by the
conditional-law definition. The two substantive branches are those
of the source's Figure~3.
\end{hbmethod}

\begin{hbtheorem}{thm:p16-idc-termination-soundness}{proved=annex:p16.idc-termination-soundness}{Termination and soundness of every successful return}
IDC performs at most \(|Z|\) exchanges and one terminal ID call.
It therefore terminates. Every successful return is the original
\(P_x(y\mid z)\), with a positive terminal denominator.
Its retained trace is a finite derivation by Rule 2, the already
proved ID do-calculus derivations, and ordinary normalization.
It can identify a conditional effect even when the original joint
\(P_x(y,z)\) is not identified.
\end{hbtheorem}

\subsection{The conditional-versus-joint necessity argument}

Put \(A=X\cup E\) and \(W=Z\setminus E\) after maximal
exchange. The source's Theorem~6 asserts that the original conditional
query is identifiable exactly when \(P_a(y,w)\) is identifiable.
Its sufficient direction is the normalization just proved. Its necessary
direction requires a conditional countermodel argument, rather than the
normalization of an unidentified joint. The following cases give complete
necessity proofs with explicit hypotheses.

\begin{hbtheorem}{thm:p16-idc-known-denominator}{proved=annex:p16.idc-known-denominator}{Conditional necessity when the terminal conditioning marginal is identified}
Suppose the terminal marginal \(P_a(W)\) is graph-uniformly
identified on positive laws. Then the original conditional query is
graph-uniformly identified if and only if \(P_a(Y,W)\) is
graph-uniformly identified. In this case IDC is complete: failure
of the terminal ID call supplies two positive models with different
original conditional effects. The empty \(W\) case is included.
\end{hbtheorem}

\begin{hbtheorem}{thm:p16-idc-response-rooted-failure}{proved=annex:p16.idc-response-rooted-failure}{A conditional countermodel when the hedge roots are response vertices}
Suppose a hedge for the terminal joint query \((A,Y\cup W)\)
has common roots \(K\subseteq Y\). Then the original conditional
query is not graph-uniformly identified on positive laws. The positive
binary forest models, extended by independent outside coordinates,
already give such conditional countermodels. No identification of
the marginal \(P_a(W)\) is required for this conclusion.
\end{hbtheorem}

\begin{hbtheorem}{thm:p16-idc-parent-root-gadget}{proved=annex:p16.idc-parent-root-gadget}{A conditional countermodel from an external parent of a conditioning root}
Suppose a hedge for the terminal joint query has every common root
in \(W\), and has a selected root \(r\) with an allowed
parent \(s\in Y\) outside its larger forest. Then two positive
binary models on the full declared graph have equal observed laws
and different conditional effects on \(Y\) given \(W\).
One may take \(s\) to be an independent fair bit and add it
to the parity equation at \(r\) in both forest models.
\end{hbtheorem}

\begin{hbtheorem}{thm:p16-idc-hedge-scope-counterexample}{proved=annex:p16.idc-hedge-scope-counterexample}{Extra nonaction vertices can belong to the larger hedge}
Take directed edges \(L\to X\to W\) and bidirected edges
\(X\leftrightarrow L\leftrightarrow W\). The whole selected
graph \(F\) and the singleton \(F'=\{W\}\) form a hedge
for action \(\{X\}\), response \(\{W\}\), but
\(F\not\subseteq F'\cup\{X\}\). Thus that inclusion
cannot be used as a general hedge lemma.

For a positive realization, take independent fair bits \(U,V\)
and private \(E\) with one-probability \(1/4\). Use
\(L=U\oplus V\), \(X=L\oplus U\), and either
\(W=X\oplus V\oplus E\) or \(W=E\).
Both observed laws have \(L,X\) independent fair and
\(W\) independent with one-probability \(1/4\), but their
response-one effects under either \(x\) are \(1/2\) and \(1/4\).
\end{hbtheorem}

\begin{hblemma}{thm:p16-idc-positive-path-gadget}{proved=annex:p16.idc-positive-path-gadget}{A finite positive gadget for an active path}
Let \(w\) be a sink in an observed/latent-root DAG, and suppose
it is d-connected to an observed set \(Y\) given an observed
set \(C\), all disjoint. There is a finite vector-valued model
on the allowed graph with a strictly positive observed law, a single
bit \(K\) as the output at \(w\), a bit coordinate \(B\)
at some vertex of \(Y\), and an assignment \(c\) such that
the two-by-two conditional table of \(K,B\) given \(C=c\)
has nonzero determinant. The bit \(K\) is not used by another
equation. Extra allowed edges may be ignored.
\end{hblemma}

\begin{hblemma}{thm:p16-idc-unused-conditioning-bit}{proved=annex:p16.idc-unused-conditioning-bit}{A conditioning-marginal witness with an unused differing coordinate}
If \(P_a(W)\) is not graph-uniformly identified on positive
laws, two finite positive models with the same observed law can
be chosen so that some bit coordinate \(L\) of a vertex
\(w\in W\) is unused in every other equation, and the
conditional probability of \(L=1\) given all the other observed
coordinates at vertices of \(W\) differs between the two models
at the same action and remaining-coordinate assignment.
\end{hblemma}

\begin{hbtheorem}{thm:p16-idc-terminal-reduction}{proved=annex:p16.idc-terminal-reduction}{The complete terminal conditional-versus-joint reduction}
After maximal graphical exchange, let \(A=X\cup E\) and
\(W=Z\setminus E\). On the graph-uniform positive-law class,
the original conditional query is identified if and only if
\(P_a(Y,W)\) is identified. If the terminal joint is not
identified, there are positive finite models with the same full
observed law and different original conditional effects. The proof
handles both an identified and an unidentified conditioning marginal.
It uses finite additional coordinates, so it makes no assertion of
completeness for a separately prescribed small alphabet family.
\end{hbtheorem}

\begin{hbcorollary}{thm:p16-idc-completeness}{proved=annex:p16.idc-completeness}{IDC, the terminal hedge criterion and do-calculus completeness}
IDC is sound and complete for graph-uniform conditional identification
on positive laws. Its result is independent of exchange order.
A query is identified exactly when there is no hedge for the full
terminal query \((X\cup E,Y\cup(Z\setminus E))\).
Every identified conditional query has a finite do-calculus derivation
with ordinary probability operations. The terminal failure record and
the conditional-witness construction together justify nonidentification;
failure of an arbitrary original joint alone would not do so.
\end{hbcorollary}

The general construction retains the two cases of the source's Theorem~6
argument. Its active-path transmission is made explicit by separate
coordinate channels and positive noise. It does not use the general
hedge inclusion rejected above, which appears as the source's Appendix
Lemma~5. The larger finite alphabets are part of the graph-uniform
countermodel construction, not a change to a prescribed input family.

\subsection{Complete calculated runs}

\begin{hbtheorem}{thm:p16-idc-worked}{proved=annex:p16.idc-worked}{Conditioning helps in a chain and obstructs identification at a collider}
For \(X\to Z\to Y,X\leftrightarrow Z\), IDC exchanges
\(Z\) for an action and returns
\(P_x(y\mid z)=p(y\mid z)\). With independent
\(U,E,E_Y\), \(U\) fair and \(E,E_Y\) of one-probability
\(1/4\), compare \(X=U,Z=X\oplus U\oplus E,Y=Z\oplus E_Y\)
with the model replacing \(Z\) by \(E\). Their common table,
for each \(x\), is
\[
\begin{array}{c|cc}
z&p(x,z,0)&p(x,z,1)\\\hline
0&9/32&3/32\\
1&1/32&3/32
\end{array}
\]
The original joint intervention law on \(Z,Y\) differs between
the models, but the conditional response-one probabilities are
\(1/4\) at \(z=0\) and \(3/4\) at \(z=1\) in both.

For \(X\to Z\leftarrow Y,X\leftrightarrow Z\), let
\(U,V\) be independent fair bits and \(E\) an independent
bit of one-probability \(1/4\). Use \(X=U,Y=V\) and either
\(Z=X\oplus U\oplus Y\oplus E\) or \(Z=Y\oplus E\).
The common observed law is \(3/16\) when \(z=y\) and
\(1/16\) when \(z\ne y\), for each \(x,y,z\).
The marginal effect on \(Y\) is the fair law in both models,
but the conditional response-one values are
\[
\begin{gathered}
 P^1_x(Y=1\mid Z=z)=1/2,\\
 P^2_x(Y=1\mid Z=0)=1/4,\\
 P^2_x(Y=1\mid Z=1)=3/4.
\end{gathered}
\]
No exchange is licensed, and the terminal joint ID call fails.
The annex gives both complete graph traces and calculations.
\end{hbtheorem}

\begin{proofboundary}{def:p16-idc-support-boundary}{Nonpositive evidence and distinct source versions}
At a nonpositive observed law, IDC's Rule 2 step requires the actual
current conditioning mass positive; terminal normalization requires the
actual terminal conditioning marginal positive. A terminal joint law
returned by the support-aware procedures of the preceding section can
be normalized only after these requirements are proved. The original
action's observed support alone does not supply them.

The UAI version is the eight-page author-hosted R-329 text dated May 2006.
Its earlier February submission, its 2007 addendum, and the later arXiv
deposit are distinct records. The addendum's correction concerns the
separate Tian conditional algorithm. Its original four-phase
cancellation procedure, full soundness and positive graph-uniform
completeness are supplied in
\cref{thm:p16-tcc-cancellation-soundness,thm:p16-tcc-completeness}.
The corrected original comparison and its complete countermodel
calculation are \cref{thm:p16-tcc-corrected-comparison}
\cite{Tian2004ConditionalCausalEffects,Shpitser2007ConditionalIdentificationAppendum}.
\end{proofboundary}
